\section{Formal Problem Definition}

\subsection{Introduction}

We consider an online variant of the Temporal Hitting Set problem under a sliding window model. 
To this end, we model the input as a temporal hypergraph $H = (V, E)$, where each hyperedge $(S, t)$ consist of a subset of vertices $S \subseteq V$ and a discrete time step $t \in \mathbb{T}$.

Given a time window of size $\epsilon$, we define the set of active hyperedges at time $t$ as:
\[
W(t) = \{ (S, t') \in E \mid t - \epsilon \le t' \le t \}.
\]

At each time step $t$, the algorithm maintains a hitting set $X(t) \subseteq V$ such that:
\[
\forall (S, t') \in W(t), \quad X(t) \cap S \neq \emptyset.
\]

\subsection{Objective}

The goal is to maintain a feasible hitting set over time while optimizing the following criteria:

\begin{itemize}
    \item \textbf{Solution size:} minimize $|X(t)|$ at each time step;
    \item \textbf{Update cost:} minimize the number of new nodes taken in consideration when updating the hitting set from $X(t-1)$ to $X(t)$;
    \item \textbf{Computational efficiency:} minimize the time required to update the solution as the sliding window evolves.
\end{itemize}

Overall, the objective is to maintain a small and stable hitting set while efficiently handling the insertion and expiration of hyperedges induced by the sliding window.